\subsection*{Counting sort}

Пусть известно, что множество объектов, массив из которых нам придется сортировать, невелико по размеру.
Тогда можно посчитать, сколько раз встречается каждый элемент, а потом вывести его нужное число раз.
Пусть множество сортируемых элементов имеет размер K, тогда время равно $O(N + K)$ и доппамять $O(K)$

\subsection*{Stable counting sort}

Пусть имеется массив A длины N -- исходный и массив B -- куда будем писать ответ. Также заведем массив P размера K , где K -- размер множества элементов.

\begin{enumerate}
    \item Пройдем по массиву A и запишем в P[i] число объектов с ключом i.
    \item Посчитаем префиксные суммы массива P, тогда мы знаем, начиная с какого индекса массива B надо писать структуру с ключом i.
    \item Идем по массиву A и пишем в нужное место, поддерживая, куда писать следующий элемент с ключом k.
\end{enumerate}

\subsection*{Least Significant Digit (LSD)}

Алгоритм Least Significant Digit сортирует целые беззнаковые числа за
линейное время.
unsigned int хранятся в виде строки из четырех байтов.
Тогда будем сортировать побайтово данные числа как
двоичные строки, сначала сортируем по убыванию последнего байта,
потом стабильно по убыванию второго байта и так далее. Важно использовать стабильную сортировку (к примеру сортировку подсчетом). После такого
получим, что массив чисел отсортирован.
Работает это за $O(Nk)$ времени, где k это число байтов или 4, то есть
за $O(N)$.

\pagebreak